\documentclass[a4paper, 12pt]{article}
\usepackage[german]{babel} % deutsch, deutsche Rechtschreibung
\usepackage[utf8]{inputenc} % Unicode-Zeichensatz als Text-Quelle
\usepackage[T1]{fontenc} % Umlaute und deutsches Trennen
\usepackage{mathptmx} % Times New Roman, gewohnter Font
%\usepackage{uarial} % Setzt den gesamten Text in Arial
\usepackage{courier} % einen schickeren Schreibmaschinenfont
\usepackage[scaled=.95]{helvet} % was serifenloses, wenn gebraucht
\usepackage{graphicx} % wir wollen Bilder einfügen
\usepackage{xfrac} % schöne Brüche im Fließtext mit sfrac

\usepackage{hyperref}
  
\usepackage{listings} % Schöne Quellcode-Listings [minted wäre besser]


\lstset{basicstyle=\sffamily, columns=[l]flexible, mathescape=true, 
  showstringspaces=false, numbers=left, numberstyle=\tiny}
\lstset{language=python} % und nur schöne Programmiersprachen ;-)
% und eine eigene Umgebung für Listings

\usepackage{float} % eigene Fließobjekte, kommen an beliebigen Stellen vor
\newfloat{listing}{htbp}{scl}[section] % Nummeriere je Abschnitt
\floatname{listing}{Listing} % listing ist ein Fließobjekt

% Auch wenn es anrüchig ist, man kann den Platz etwas mehr ausnützen
\usepackage[paper=a4paper,width=15cm,left=35mm,height=22cm]{geometry}
\usepackage{setspace}
\linespread{1.2} % nicht ganz anderthalbzeilig, nur ein bisschen mehr Platz - war 1.15
\setlength{\parskip}{0.4em} % kleiner Paragraphen(Absatz)-abstand
\setlength{\parindent}{0em} % im Deutschen Einrückung nicht üblich

% Seitenmarkierungen 
\usepackage{fancyhdr} % Schickere Header und Footer
\pagestyle{fancy}
% Zeichensatz für Header/Footer
\newcommand{\phv}{\fontfamily{phv}\fontseries{m}\fontsize{9}{11}\selectfont}
\fancyhead[L]{\phv \leftmark} % Kurztitel links oben
\fancyhead[R]{\phv \thepage} % rechts oben die Seitenzahl
\fancyfoot[L]{\phv Hochschule Mannheim} % Institution links unten
\fancyfoot[C]{\ } % keine Seitenzahl unten Mitte
\fancyfoot[R]{\phv Medizintechnik} % Studiengang rechts unten

\usepackage{url} % wir wollen eine URL anzeigen


\usepackage{titlesec}

\begin{document}

\begin{titlepage}
	\centering
	%\vspace*{2cm} %Abstand nach oben

	% Logos:	
    \begin{minipage}{0.1\textwidth}
        \includegraphics[height=2.5cm]
        {Hochschule_Mannheim_logo.png}
    \end{minipage}
    \hfill 
    \begin{minipage}{0.25\textwidth}
        \includegraphics[height=5cm]
        {loewenstein_logo.png}
    \end{minipage}	

    \vspace{2.5cm} % Abstand nach den Bildern
	
	%Titel:
	{\Huge\bfseries Praxissemesterbericht \par}
    \vspace{2cm} % Abstand nach dem Titel

    % titlepage content
    \begin{tabular}{ r  l }
        \textbf{Autor:} & Rebekka Hahn \\
        \textbf{Matrikelnummer:} & 1921861 \\
        \textbf{Semester:} & 11. Semester \\
        \textbf{Studiengang:} & Medizintechnik \\
        \textbf{Beginn Praxissemester:} & 02.09.2024 \\
        \textbf{Ende Praxissemester:} & 28.02.2025 \\
        \textbf{Firma:} & Löwenstein medical \\
        \textbf{Betreuer:} & Patrick von Poblotzki, Christoph Elsner \\
    \end{tabular}
    
    \vfill % Abstand nach unten

    % Ort und Datum
    {\large Ludwigshafen am Rhein, \today \par}

\end{titlepage}

\newpage
\tableofcontents            % Inhaltsverzeichnis generieren
\newpage                    % Neue Seite für den Textteil

{\bfseries Selbständigkeitserklärung}\\ \\
Ich versichere, dass ich diesen PS-Bericht selbständig und nur unter Verwendung der angegebenen
Quellen und Hilfsmittel angefertigt habe. Die Stellen, an denen Inhalte aus den Quellen verwendet
wurden, sind als solche eindeutig gekennzeichnet. Die Arbeit hat in gleicher oder ähnlicher Form bei
keinem anderen Prüfungsverfahren vorgelegen. \\
\vspace{1.0cm} \\
\line(1,0){430} \\
Datum, Ort und Unterschrift\\

\newpage
\section{Einleitung}\label{Einleitung} 
Die Einleitung zu meinem Praxissemesterbericht
Kann ja was darüber erzählen, wo ich bin, was Abteilung, Standort

\newpage
\section{Schulungen}\label{Schulungen}
Schulungen vom Betrieb

\subsection{Arbeitssicherheit}\label{Arbeitssicherheit}
Was muss man im Büro bei einem Notfall tun, bzw sich an wen wenden?

\subsection{IT-Sicherheit}\label{IT-Sicherheit}
Richtiger Umgang mit firmeninternen, sensiblen Daten

\subsubsection{Phishing}\label{Phishing}
phishing mails

\subsubsection{DSGVO}\label{DSGVO}
DSGVO \cite{dsgvo2024}

\subsection{ESD}\label{ESD}
= Electrostatic Discharge

\subsubsection{Physik ESD}\label{physik_esd}
elektrostatik

\subsubsection{Maßnahmen}\label{measurements_esd}
Zum Schutz der Hardware vor elektrostatischer Entladung (ESD) werden geerdete Arbeitsmatten und Handgelenk-Erdungsbänder verwendet. Für den sicheren Transport zwischen Arbeitsplätzen wird die Hardware in speziellen, abschirmenden Verpackungen transportiert, die wie ein Faradayscher Käfig wirken und so elektrostatische Felder abschirmen.

\subsection{Dokumentation}\label{Dokumentation}
Dokumente sind Unterlagen, die nach festgelegten Prozessen verwaltet und regelmäßigen Revisionen unterzogen werden. \\
\subsubsection{Lenkung Dokumente}\label{lenkung_dokumente}
Angemessenheit

\subsection{Qualitätsmanagement}\label{Qualitätsmanagement}
MDR, IVDR, AIMDD, etc


\subsection{Meetings}\label{Meetings}
SCRUM - Erklären

\subsubsection{SCRUM}\label{SCRUM}
Scrum ist ein systematischer Ansatz um Projekte strukturiert durchzuführen. Es soll die Teams bei der Lösung komplexer Probleme unterstützen indem Rollen, Regeln und Ereignisse definiert werden. Die zugrundeliegenden Prinzipien sind Empirie und Lean Thinking. 
„Empirie, die Erfahrung selbst und die auf Erfahrung beruhende Erkenntnis.  
 \cite{dorsch_empirie}"
 
Lean Thinking https://www.lean.org/lexicon-terms/lean-thinking-and-practice/
bei ruhigerer Lage mal analysieren % TODO
 

\cite{scrum2020}
% https://scrumguides.org/docs/scrumguide/v2020/2020-Scrum-Guide-German.pdf


\newpage
\section{Aufgaben}\label{Aufgaben}
Grob was ich an Aufgaben im PS hatte, aufgrund von der wirtschaftlichen Relevanz der Daten wird nur grob auf die Inhalte eingegangen

\subsection{Dokumentationsautomatisierung}\label{Dokumentationsautomatisierung}
python woop woop
\subsubsection{Regular Expression}\label{regularExpression}
python library und syntax == cool aber aufwändig \\
import re

\subsection{Library Adapter}\label{LibraryAdapter}
Eine Library durch eine aktuellere austauschen in C++

\newpage
\section{Ergebnisse}
Präsentiere und diskutiere hier die Ergebnisse deines Berichts.

\newpage
\section{Fazit}
Im Fazit fasst du alles zusammen und gibst einen Ausblick.

%  Literaturverzeichnis
\newpage
\bibliographystyle{unsrt} %unsrt
\bibliography{Praxissemesterbericht_bibliographie}

\end{document}
